\documentclass[12pt]{article}
\usepackage[utf8]{inputenc}
\usepackage[T1]{fontenc}
\usepackage{hyperref}
\title{Samay Language Resources for K-12 Curriculum Development}
\author{Compiled by Claude Sonnet 5.5 (claude-sonnet-5-5) for Andr\'as Kornai}
\date{}
\begin{document}
\maketitle

\section{Introduction}
Samay (Osamayi, Sama, Samaye, Shamay) is a minor Bantu language of northeastern Gabon. Ethnologue gave about 1,000 speakers in 2007, according to Wikipedia. Little is documented. French and larger regional languages are reported to dominate, especially among younger people.

\section{Primary Educational Resources}
No primers, textbooks or school programmes for Samay were found.

\section{Linguistic and Reference Materials}
\begin{itemize}
\item Wikipedia, Samay language: a stub with basic identification. \url{https://en.wikipedia.org/wiki/Samay_language}
\item Glottolog entry for Osamayi (osam1235), with classification. \url{https://glottolog.org/resource/languoid/id/osam1235.bigmap.html}
\end{itemize}
No dictionary or grammar was found.

\section{Religious and Community Resources}
\begin{itemize}
\item Scripture Earth page for syx, which lists scripture resources. \url{https://www.scriptureearth.org/00eng.php?iso=syx}
\item Joshua Project profile of the Samay people in Gabon. \url{https://joshuaproject.net/people_groups/20995}
\end{itemize}
It is unclear what scripture material exists. The sources conflict on this.

\section{Digital and Audio Resources}
\begin{itemize}
\item Global Recordings Network page for syx. It states that no recordings are currently available. \url{https://globalrecordings.net/en/language/syx}
\end{itemize}
Scripture Earth mentions audio through Bible.is and Global Recordings Network, which conflicts with the Global Recordings page. This could not be resolved.

\section{Curriculum Development Context}
Almost nothing exists for Samay. A curriculum would first need fieldwork, an orthography and community consultation to establish how many young people still speak the language. The French-medium Gabonese school system is the likely context. Any native-language component would have to be built from scratch.
\end{document}